\chapter{前言}\label{preamble}

\topicindex{nodes}
\topicindex{boxes}
\topicindex{\LUA}

这是一本参考手册，不是教程。这意味着我们讨论的是相对于传统 \TEX\ 的变
化，同时也介绍新功能。因此，我们会引用一些假定读者已知的概念，或者将在
后面解释的概念。

一般用户不需要了解本手册中的太多内容。例如，字体和语言通常由你所使用的
宏包来处理。在用户层面，通常也不需要真正去摆弄节点列表。如果你确实要摆
弄，那最好清楚自己在做什么。阅读 Donald Knuth 的 ``The
\TEX\ Book'' 是很好的时间投资，因为了解一切的起源也是有益的。
Victor Eijkhout 的 ``\TEX\ by Topic'' 提供了更具概括性的
综述。你也可以翻阅 ``\ETEX\ 手册'' 以及关于 \PDFTEX\ 的
文档。

但是，如果你是因为 \LUA\ 来到这里的，那么你只需要知道：可以在排版运行
中调用它。你使用的宏包可能会提供一些封装机制，但完成这项工作的基本
\lpr {directlua} 命令是：

\begin{typing}
\directlua{tex.print("Hi there")}
\end{typing}

你可以将代码放在花括号之间，但如果代码量较大，也可以将其放入文件中，并
使用常规的 \LUA\ 命令加载该文件。

如果你仍然决定继续阅读，那么了解什么是节点会很有帮助，所以我们在这里做
一个简要介绍。如果你输入以下这段文字：

\begin{typing}
Hi There
\end{typing}

最终我们会得到一个节点链表，用 \ASCII\ 艺术表示如下：

\begin{typing}
H <=> i <=> [glue] <=> T <=> h <=> e <=> r <=> e
\end{typing}

当我们有一个段落时，实际上会得到这样的结构：

\begin{typing}
[localpar] <=> H <=> i <=> [glue] <=> T <=> h <=> e <=> r <=> e <=> [glue]
\end{typing}

每个字符都会成为一个所谓的 glyph 节点，这是一个包含当前字体、字符编码
和当前语言等属性的记录。空格会成为 glue 节点。节点类型有很多，我们将在
后面讨论。每个节点指向前一个节点或后一个节点（如果它们存在的话）。

事先了解以下几点也是有益的：\TEX\ 基本上围绕着创建段落和页面来运作。
段落构建器接收一个列表并将其断行为多行。我们将水平材质转换为垂直材质。
行就是所谓的 box，可以用 glue、penalty 等分隔。页面构建器累积行，在合
适时触发输出例程，该例程会接收目前为止的列表。构建实际页面并不是
\TEX\ 的职责，而是通过允许操作 box 的原始命令来完成的。结果交还给
\TEX{}，并输出到（通常是 \PDF）文件。

\LUATEX\ 引擎在处理过程中几乎每个合理的节点都为 \LUA\ 代码提供了钩子：
收集内容、断字、应用字体特性、断行等。这意味着你可以重载 \TEX\ 的固有
行为，而这种固有行为仍然是基准。当我们提到 ``回调（callbacks）'' 时，
指的就是这些钩子。

纯 \TEX\ 基本上是一个用于编写特定样式的基础框架，而像 \CONTEXT\ 和
\LATEX\ 这样的宏包则为用户提供了大量额外工具，使文档排版美观。它们隐
藏了字体管理、语言需求、将结构转化为排版结果、分页、包含图像等繁琐细
节。你应该注意，当你接入自己的代码来操作列表时，这可能会与你所使用的宏
包产生冲突。

在接下来的章节中阅读关于节点的内容时，最好记住与它们相关的命令。以下
是几个例子：

\begin{longtable}{l l p{5cm}}
\toprule 命令                & 节点          & 说明 \\
\midrule
\prm {hbox}            & \nod {hlist} & 水平 box \\
\prm {vbox}            & \nod {vlist} & 垂直 box，基线位于底部 \\
\prm {vtop}            & \nod {vlist} & 垂直 box，基线位于顶部 \\
\prm {hskip}           & \nod {glue}  & 水平间距，可伸缩可收缩 \\
\prm {vskip}           & \nod {glue}  & 垂直间距，可伸缩可收缩 \\
\prm {kern}            & \nod {kern}  & 水平或垂直固定间距 \\
\prm {discretionary}   & \nod {disc}  & 断字点（前置、后置、替代） \\
\prm {char}            & \nod {glyph} & 一个字符 \\
\prm {hrule}           & \nod {rule}  & 一条水平线 \\
\prm {vrule}           & \nod {rule}  & 一条垂直线 \\
\prm {textdir(ection)} & \nod {dir}   & 文本方向的改变 \\
\bottomrule
\end{longtable}

目前这些内容应该足以让你理解接下来的章节了。
